% Irodov, user-supplied book, section 2.1 original problems 12--20,
% printed p.64 / PDF p.73.
\thermoproblem{12}{64}
一摩尔理想气体沿下列过程变化：
$T=T_0+\alpha V^2$，
其中 $T_0,\alpha$ 为正常数，
$V$ 为该摩尔气体的体积。
求过程中能够达到的最小压强，
并定性画出 $p$ 随 $V$ 变化的图线。
\exsolution
一摩尔理想气体满足 $pV=RT$。
代入过程方程：
\[
p(V)=\frac{R(T_0+\alpha V^2)}{V}
=R\left(\frac{T_0}{V}+\alpha V\right),
\quad V>0.
\]
对体积求导：
\[
\frac{\mathrm dp}{\mathrm dV}
=R\left(-\frac{T_0}{V^2}+\alpha\right).
\]
极小值在
$V_*=\sqrt{T_0/\alpha}$，
此点二阶导数
$2RT_0/V_*^3>0$。
于是
\[
\boxed{p_{\min}=p(V_*)
=2R\sqrt{\alpha T_0}}.
\]
图线在 $V\to0^+$ 时由 $RT_0/V$
主导而趋无穷，在 $V\to\infty$ 时
由 $R\alpha V$ 主导亦趋无穷，
在 $V_*$ 处有唯一谷底。
\excheck 过程曲线为 U 形，
不是理想气体等温过程的双曲线；
这里气体温度随体积而变化。

\thermoproblem{13}{64}
在均匀重力场中放置一个高圆柱形容器，
内装氮气。
自由落体加速度大小为 $g$，
气体沿高度 $h$ 的温度分布使得
其\emph{质量密度处处相同}。
求温度梯度 $\mathrm dT/\mathrm dh$。
\exsolution
高度向上取正，
静力学平衡给出
$\mathrm dp/\mathrm dh=-\rho g$。
由理想气体状态方程
$p=\rho RT/M$，
其中氮气摩尔质量 $M$、
$\rho$ 为质量密度。
题设 $\rho$ 在所有高度都相同，
所以
\[
\frac{\mathrm dp}{\mathrm dh}
=\frac{\rho R}{M}\frac{\mathrm dT}{\mathrm dh}.
\]
与静力方程联立：
\[
\boxed{\frac{\mathrm dT}{\mathrm dh}
=-\frac{Mg}{R}}.
\]
以氮气 $M=0.0280\,\mathrm{kg/mol}$
及 $g=9.81\,\mathrm{m/s^2}$，
得
$\boxed{\mathrm dT/\mathrm dh
\simeq-0.033\,\mathrm{K/m}=-33\,\mathrm{mK/m}}$。
\excheck 温度随高度线性下降，
与常温气体密度不随高度变化这一特定假设相容；
该梯度不是任意大气环境的实际气温直减率。

\thermoproblem{14}{64}
假设大气中的压强与质量密度满足
$p/\rho^{\,n}=\text{常数}$，
其中 $n$ 与高度无关。
利用静力学平衡和理想气体状态方程，
求相应温度梯度。
\exsolution
写多方关系为 $p=C\rho^n$，
则气体温度满足
\[
T=\frac{Mp}{R\rho}
=\frac{MC}{R}\rho^{n-1}.
\]
两边求导：
\[
\frac{\mathrm dT}{T}
=(n-1)\frac{\mathrm d\rho}{\rho},
\qquad
\frac{\mathrm dp}{p}=n\frac{\mathrm d\rho}{\rho}.
\]
由静力方程
$\mathrm dp/\mathrm dh=-\rho g$，
以及 $p/\rho=RT/M$，
得到
\[
\frac{\mathrm dT}{\mathrm dh}
=(n-1)\frac{T}{n p}\frac{\mathrm dp}{\mathrm dh}
=-(n-1)\frac{T\rho g}{np}.
\]
化简：
\[
\boxed{\frac{\mathrm dT}{\mathrm dh}
=-\frac{Mg}{nR}(n-1)}.
\]
\excheck $n=1$ 时等温且温度梯度为零；
如果 $n>1$，温度随高度升高而降低。

\thermoproblem{15}{64}
地表附近空气处于标准状况，
假设温度和平均摩尔质量随高度不变。
求地表上方 $5.0\,\mathrm{km}$ 处，
以及地面下方深 $5.0\,\mathrm{km}$ 的矿井中的气压。
\exsolution
等温静力平衡的气压分布为
\[
p(h)=p_0\exp\left(-\frac{Mgh}{RT}\right),
\]
其中高度 $h$ 在地面以上取正、矿井中取负。
取 $p_0=1.00\,\mathrm{atm}$、
$T=273.15\,\mathrm K$、
$M\simeq0.029\,\mathrm{kg/mol}$。
定义
$a=M g(5000\,\mathrm m)/(RT)\simeq0.63$。
于是
\[
\boxed{p(+5.0\,\mathrm{km})
=p_0e^{-a}\simeq0.53\,\mathrm{atm}\approx0.5\,\mathrm{atm}},
\]
\[
\boxed{p(-5.0\,\mathrm{km})
=p_0e^{a}\simeq1.9\,\mathrm{atm}\approx2\,\mathrm{atm}}.
\]
\excheck 结果来自温度在上下5 km范围内均恒定的
理想化模型；现实矿井的地温梯度与空气流动会改变数值。

\thermoproblem{16}{64}
在均匀重力场中，
假设空气温度为 $0^\circ\mathrm C$、
平均摩尔质量 $M$ 与 $g$ 均不随高度改变。
求两个高度之间的距离 $\Delta h$，
使相应空气密度分别：
(a) 相差 $e$ 倍；
(b) 相差 $\eta=1.0\%$。
\exsolution
等温大气的密度与气压同指数衰减：
\[
\rho(h)=\rho(0)e^{-h/H},
\qquad H=\frac{RT}{Mg}
\]
是密度标高。
两点间距 $\Delta h>0$ 的密度比
$\rho(h)/\rho(h+\Delta h)=e^{\Delta h/H}$。
(a) 若密度相差 $e$ 倍，
$\Delta h/H=\ln e=1$，所以
\[
\boxed{\Delta h_a=\frac{RT}{Mg}}.
\]
(b) 若相差 $1+\eta$ 倍，
\[
\boxed{\Delta h_b=\frac{RT}{Mg}\ln(1+\eta)
\simeq\eta\frac{RT}{Mg}}.
\]
取 $T=273.15\,\mathrm K$、
$M=0.029\,\mathrm{kg/mol}$，
得到
$\boxed{\Delta h_a\simeq8.0\,\mathrm{km}}$、
$\boxed{\Delta h_b\simeq0.080\,\mathrm{km}=80\,\mathrm m}$。
\excheck 当 $\eta\ll1$ 时
$\ln(1+\eta)\simeq\eta$；
不能把 $1.0\%$ 错当成倍率 $1.0$。

\thermoproblem{17}{64}
一个高为 $h$、底面积为 $S$ 的竖直圆柱容器中
装有摩尔质量为 $M$ 的等温理想气体，
温度为 $T$。
容器底面压强为 $p_0$，
重力加速度为常数 $g$。
求容器内气体总质量 $m$。
\exsolution
底端 $z=0$ 向上取高度 $z$。
等温静力平衡给出
\[
p(z)=p_0\exp(-Mg z/(RT)).
\]
理想气体质量密度为
$\rho(z)=Mp(z)/(RT)$。
因此总质量
\[
m=S\int_0^h\rho(z)\,\mathrm dz
=\frac{MS p_0}{RT}
\int_0^h\exp(-Mg z/(RT))\,\mathrm dz.
\]
积分得
\[
\boxed{m=\frac{S p_0}{g}
\left(1-e^{-Mgh/(RT)}\right)}.
\]
\excheck 对薄层高度 $h\ll RT/(Mg)$，
指数展开给
$m\simeq S p_0 M h/(RT)$，
恢复均匀近似密度 $\rho_0$ 乘体积 $Sh$。

\thermoproblem{18}{64}
把摩尔质量为 $M$ 的等温理想气体
放在均匀重力场中一个\emph{非常高}的竖直容器里，
温度为 $T$、重力加速度为 $g$。
求气体质量分布的质心高度
$h_C$（相对容器底面）。
\exsolution
等温大气密度分布为
$\rho(h)=\rho_0e^{-h/H}$，
其中 $H=RT/(Mg)$。
容器足够高，可用上限无穷大。
质心高度是质量加权平均：
\[
h_C=\frac{\int_0^\infty h\,\rho(h)S\,\mathrm dh}
{\int_0^\infty\rho(h)S\,\mathrm dh}
=\frac{\int_0^\infty h e^{-h/H}\,\mathrm dh}
{\int_0^\infty e^{-h/H}\,\mathrm dh}.
\]
由
$\int_0^\infty e^{-h/H}\mathrm dh=H$、
$\int_0^\infty h e^{-h/H}\mathrm dh=H^2$，
得到
\[
\boxed{h_C=H=\frac{RT}{Mg}}.
\]
\excheck 质心高度不受容器截面积与底部压强影响，
是这一指数密度模型的标高。

\thermoproblem{19}{64}
在均匀重力场中，理想气体摩尔质量为 $M$、
自由落体加速度为 $g$。
已知 $h=0$ 时压强 $p=p_0$，
若温度随高度分别满足
(a) $T=T_0(1-ah)$；
(b) $T=T_0(1+ah)$，
其中 $a>0$，
求相应压强分布 $p(h)$。
\exsolution
静力学平衡为
$\mathrm dp/\mathrm dh=-\rho g$，
用
$\rho=Mp/(RT(h))$ 消去密度：
\[
\frac{\mathrm dp}{p}
=-\frac{Mg}{RT(h)}\,\mathrm dh.
\]
记
$\nu=Mg/(aRT_0)>0$。
(a) 对 $T(h)=T_0(1-ah)$，
\[
\ln\frac p{p_0}
=-\frac{Mg}{RT_0}\int_0^h\frac{\mathrm dz}{1-az}
=\nu\ln(1-ah),
\]
所以
\[
\boxed{p_a(h)=p_0(1-ah)^\nu},
\qquad h<1/a.
\]
(b) 对 $T(h)=T_0(1+ah)$，
\[
\ln\frac p{p_0}
=-\frac{Mg}{RT_0}\int_0^h\frac{\mathrm dz}{1+az}
=-\nu\ln(1+ah),
\]
所以
\[
\boxed{p_b(h)=p_0(1+ah)^{-\nu}}.
\]
\excheck 两式在 $h=0$ 均给 $p_0$；
(a) 的物理定义域须保证温度 $T(h)>0$。

\thermoproblem{20}{64}
一根水平圆柱形管一端封闭、
另一端在旋转轴上与外界大气连通。
管绕穿过其开放端的竖直轴以角速度
$\omega$ 匀速旋转。
管内空气处于温度 $T$ 的等温平衡，
摩尔质量为 $M$，开放端大气压为 $p_0$。
求管内气压随距旋转轴距离 $r$ 的分布。
\exsolution
在随管转动的参考系中，
空气相对管静止。
气体在沿管轴径向满足离心静力平衡：
\[
\frac{\mathrm dp}{\mathrm dr}=\rho\omega^2r.
\]
用等温理想气体关系
$\rho=Mp/(RT)$ 代入：
\[
\frac{\mathrm dp}{p}
=\frac{M\omega^2}{RT}r\,\mathrm dr.
\]
由开放端边界 $p(0)=p_0$ 积分：
\[
\ln\frac{p(r)}{p_0}
=\frac{M\omega^2r^2}{2RT},
\]
所以
\[
\boxed{p(r)=p_0
\exp\left(\frac{M\omega^2r^2}{2RT}\right)}.
\]
\excheck 离心平衡使压强随半径\emph{增加}；
本题不是地球大气沿竖直高度衰减的模型。
